\hspace{3mm}

\centerline{\bf \Huge Тренировка 1.}

\begin{flushright}
        \scriptsize{— В каждом есть другая личность. Две личности и создают одно сознание.\\
        — Две?\\ 
        — Тот, на кого смотрят другие люди, и тот, кто сам смотрит на них.\\
        Евангелион. Эпизод 16. Смертельная болезнь.
}
\end{flushright}


\problem{1}
При каких  натуральных $n$ число $20^n + 16^n - 3^n - 1$ делится на $323$?

\problem{2}
На доске написано простое число. Каждую минуту его изменяют: умножают на 2 и либо прибавляют, либо вычитают 1. Докажите, что когда-нибудь получится составное число.

\problem{3}
Докажите, что число $2^{n!}$ даёт остаток 1 при делении на любое нечётное $m \leq n.$

\problem{4}
В треугольнике $A B C$ известно, что $\angle A=15^{\circ}$ и $\angle B=30^{\circ}$. На стороне $A B$ отметили такую точку $E$, что $\angle A C E=90^{\circ}$. Найдите отрезок $A E$, если $B C=2.$

\problem{5}
На гипотенузе $BC$ прямоугольного треугольника $ABC$ выбрана точка $K$ так, что $AB = AK$. Отрезок $AK$ пересекает биссектрису $CL$ в её середине. Найдите острые углы треугольника $ABC$.

\problem{6}
Точка $K$ --- середина гипотенузы $AB$ прямоугольного равнобедренного треугольника $ABC$. Точки $L$ и $M$ выбраны на катетах $BC$ и $AC$ соответственно так, что $BL = CM$. Докажите, что треугольник $LMK$ --- также прямоугольный равнобедренный.


\problem{7}
Существуют ли такие различные натуральные числа $a, b$ и $c$, что число $a+1/a$ равно полусумме чисел $b+1/b$ и $c+1/c?$

\problem{8}
Пусть $m,n$ натуральные числа такие, что НОД$(m, n) = 1$. Найдите количество клеток, которое пересекает отрезок соединяющий начало координат и точку с координатами $(m, n)$